% =======================================================
% 2. TERMINI
% =======================================================
\section{Termini}
 
% ------------- A -------------
\phantomsection
\addcontentsline{toc}{subsection}{A}
\subsection*{A}
\begin{itemize}[leftmargin=*]
      \item \textbf{Amministratore:} Ruolo del gruppo responsabile di assicurare
      l'efficienza di procedure, strumenti e tecnologie a supporto del
      \textit{way of working}; redige e mantiene aggiornate le Norme di
      Progetto e il Piano di Progetto; gestisce la creazione delle issue tramite Google Form; verifica,
      prima della chiusura di ogni sprint, che i tempi effettivi nel foglio di
      calcolo siano stati compilati correttamente.
      \item \textbf{Analista:} Ruolo del gruppo responsabile delle attività di
      analisi dei requisiti. Attivo dalla baseline RTB.     
      \item \textbf{Attore:} In ambito UML, rappresenta un'entità esterna (un utente umano o un altro sistema) che interagisce direttamente con il sistema software per raggiungere uno specifico obiettivo.
      \item \textbf{Analisi dei Requisiti:} Documento essenziale che fissa i requisiti che il fornitore si impegna a soddisfare in accordo con il proponente. Traduce i bisogni dell'utente descrivendo esclusivamente \textit{cosa} il sistema deve fare, senza entrare in alcun dettaglio implementativo (il \textit{come}). Costituisce la base contrattuale e il documento fondamentale per l'accesso e il superamento della revisione RTB.
      \item \textbf{API (Application Programming Interface):} Insieme di regole e strumenti che permettono a diversi sistemi software di comunicare tra loro. Nel contesto del progetto, consente l'invocazione di servizi esterni di LLM tramite chiamate strutturate.
      \item \item \textbf{Appunti (Clipboard):} Area temporanea di memoria del sistema operativo utilizzata per conservare dati copiati o tagliati dall'utente, rendendoli disponibili per operazioni di incolla. Nel progetto è utilizzata per trasferire testo da e verso l'editor.
      \item \textbf{AWS (Amazon Web Services):} Piattaforma di cloud computing che fornisce servizi on-demand come elaborazione, storage e database.
      \end{itemize}
 
% ------------- B -------------
\phantomsection
\addcontentsline{toc}{subsection}{B}
\subsection*{B}
\begin{itemize}[leftmargin=*]
      \item \textbf{Backlog:} Lista ordinata per priorità di tutte le attività, funzionalità
      e requisiti da realizzare nel corso del progetto. Nel contesto del gruppo viene
      gestita tramite GitHub Project e aggiornata ad ogni Sprint
      Planning.
      \item \textbf{Baseline:} Versione stabile e approvata di un insieme di
          documenti o artefatti, fissata in un determinato momento del progetto
          e utilizzata come riferimento ufficiale per le fasi successive. Nel
          contesto del progetto didattico le baseline principali sono: candidatura
          (CC), Requirements and Technology Baseline (RTB) e Product Baseline (PB).
      \item \textbf{Brainstorming:} Tecnica di generazione collettiva delle idee adottata
      dal gruppo come metodo decisionale principale. Si articola in tre fasi:
      raccolta libera delle proposte senza valutazione immediata, confronto
      strutturato con analisi di vantaggi e svantaggi, scelta collegiale della
      soluzione ritenuta più adeguata con registrazione nel verbale.
      \item \textbf{Branch:} Linea di sviluppo indipendente all'interno di un
          repository Git, utilizzata per lavorare su modifiche isolate senza
          influenzare il ramo principale (\texttt{main}).
      \item \textbf{Burndown Chart:} Grafico che mostra l'andamento del lavoro residuo nel
      corso di uno sprint o dell'intero progetto, confrontando il ritmo di avanzamento
      effettivo con quello pianificato. Utilizzata per monitorare il rischio di ritardo
      rispetto alle milestone.
\end{itemize}
 
% ------------- C -------------
\phantomsection
\addcontentsline{toc}{subsection}{C}
\subsection*{C}
\begin{itemize}[leftmargin=*]
      \item \textbf{Candidatura (CC):} Fase iniziale del progetto didattico, antecedente
          all'aggiudicazione del capitolato, in cui il gruppo si propone come
          fornitore producendo la Valutazione dei Capitolati e la Dichiarazione degli
          Impegni. Non è un documento ma una baseline di fase.
      \item \textbf{Cerimonie:} Insieme degli incontri strutturati previsti dal
      framework Scrum: Sprint Planning, Sprint Review, Sprint
      Retrospective e Daily Stand-up. Nel gruppo le prime tre si svolgono il martedì
      in sequenza fissa; il Daily Stand-up è gestito in modo asincrono.
      \item \textbf{Chatbot:} Sistema software che consente l'interazione con l'utente tramite linguaggio naturale. Nel progetto può fungere da mediatore tra le richieste dell'utente e le funzionalità disponibili, traducendo intenzioni espresse liberamente in azioni eseguibili.
      \item \textbf{Casi d'uso:} Tecnica di modellazione che descrive le funzionalità del sistema tramite una "visione esterna" (come percepite dall'attore), ignorandone l'implementazione. A livello tecnico è definito come un insieme di scenari (sequenze di azioni che includono uno scenario principale e possibili scenari alternativi o di errore) che condividono lo stesso scopo finale per l'utente.
      \item  \textbf{CI/CD (Continuous Integration / Continuous Deployment):} Pratica di sviluppo software che prevede l’integrazione frequente del codice in un repository condiviso e la sua automazione di build, test e distribuzione.
      \item \textbf{Ciclo Di Vita Documentale:} Sequenza di fasi che ogni documento
          percorre dalla sua creazione alla pubblicazione ufficiale. Nel contesto
          del gruppo si articola in quattro fasi: Stesura (redazione del contenuto),
          Revisione (controllo da parte del Verificatore), Pubblicazione (integrazione
          nel branch \texttt{main} tramite pull request) e Ingresso in Baseline
          (la versione diventa riferimento ufficiale).
      \item \textbf{Code Coverage:} Metrica che misura la percentuale di codice sorgente eseguita durante l’esecuzione dei test automatici.
      \item \textbf{Colli di bottiglia:} Fase o componente di un processo che rallenta l'efficienza complessiva, limitando la capacità operativa del team o del sistema.
      \item \textbf{Commit:} Operazione Git che registra un insieme di modifiche
      nel repository locale, associandovi un messaggio descrittivo. Nel flusso
      del gruppo ogni commit relativo a una issue include il riferimento
      all'identificativo della issue nel messaggio, secondo la convenzione
      \texttt{git commit -m "\#<id> <messaggio>"}, garantendo la tracciabilità
      delle modifiche rispetto alle attività pianificate.
      \item \textbf{Compito:} Obiettivo espresso dall'utente in linguaggio naturale (es.\ ``applicare il cappello critico''), che il sistema interpreta per determinare la funzionalità da eseguire.
      \item \item \textbf{Consuntivo:} Documento o sezione del Piano di Progetto
      che riporta, al termine di ogni sprint, le ore effettivamente
      impiegate da ciascun membro per ciascun ruolo, il relativo costo
      effettivo e lo scostamento rispetto al preventivo.
\end{itemize}
 
% ------------- D -------------
\phantomsection
\addcontentsline{toc}{subsection}{D}
\subsection*{D}
\begin{itemize}[leftmargin=*]
      \item \textbf{Daily Stand-up:} Cerimonia agile breve e quotidiana in cui ciascun
      membro del team condivide: (1)~cosa ha fatto, (2)~cosa farà, (3)~eventuali
      blocchi o dubbi. Nel gruppo viene gestita in modo asincrono tramite
      WhatsApp, salvo i giorni in cui è prevista una riunione
      sincrona.
      \item \textbf{Debito Tecnico:} Insieme di scelte progettuali o implementative
      sub-ottimali accumulate durante lo sviluppo per ragioni di tempo o priorità,
      che richiedono interventi correttivi futuri (refactoring) e che, se non gestite,
      rallentano la velocity degli sprint successivi.
      \item \textbf{Decision Tree (Albero Di Decisione):} Modello strutturato ad albero utilizzato per rappresentare decisioni e le loro possibili conseguenze.
      \item \textbf{Definition of Done (DoD):} Lista esplicita di criteri che un
      task deve soddisfare per essere considerato completato a tutti
      gli effetti. Si articola in due livelli: livello task (criteri tecnici e
      procedurali sul singolo deliverable) e livello sprint (criteri di processo
      sull'intero sprint). Si espande al crescere del progetto: nella fase documentale
      copre redazione, revisione e tracciabilità; nella fase di sviluppo includerà
      anche criteri su test e copertura del codice.
      \item \textbf{Diario Di Bordo:} Momento di rendicontazione periodica al committente
      in cui il gruppo espone il rapporto tra costi sostenuti e avanzamento
      conseguito, comunicando le difficoltà riscontrate. Si svolge secondo il
      calendario stabilito dai docenti.
      \item \textbf{Dichiarazione Degli Impegni:} Documento gestionale esterno redatto
          in fase di candidatura che specifica: scadenza ultima di consegna prevista,
          preventivo dei costi finali calcolato secondo il regolamento, totale delle
          ore produttive assegnate a persona (nell'intervallo 80--95 ore), distribuzione
          delle ore per ruolo, rischi attesi e impegni formali assunti verso il
          committente.
      \item \textbf{Discord:} Piattaforma di comunicazione vocale e testuale utilizzata
            dal gruppo come ambiente principale per le riunioni interne e il coordinamento
            quotidiano.
      \item \textbf{Distant Writing:} Funzionalità di generazione automatica del testo tramite prompt, in cui l’LLM produce contenuti estesi a partire da istruzioni concise fornite dall’utente.
\end{itemize}
 
% ------------- E -------------
\phantomsection
\addcontentsline{toc}{subsection}{E}
\subsection*{E}
\begin{itemize}[leftmargin=*]
      \item \textbf{Endpoint:} Punto di accesso esposto da un servizio web o da un'API, identificato da un URL specifico, tramite il quale un client può inviare richieste e ottenere risposte. Nel progetto indica l'endpoint API del servizio AI utilizzato per l'elaborazione delle richieste.
\end{itemize}
 
% ------------- F -------------
\phantomsection
\addcontentsline{toc}{subsection}{F}
\subsection*{F}
\begin{itemize}[leftmargin=*]
      \item \textbf{File System:} Componente del sistema operativo responsabile della gestione dei file locali. Nel progetto può fungere da attore secondario per operazioni di upload e download delle note.
      \item \textbf{Foglio Di Calcolo:} Strumento condiviso del gruppo utilizzato per
      il monitoraggio aggregato delle misurazioni. Raccoglie i dati delle task
      registrati tramite Google Form e calcola automaticamente, per ogni sprint,
      preventivo e consuntivo di ore e costi per membro e per ruolo, oltre alle
      risorse residue. Le formule applicate sono: $\text{ore} = \text{minuti} / 60$
      e $\text{costo} = \text{ore} \times \text{costo orario del ruolo}$.
      \item \textbf{Formattazione:} Insieme delle funzionalità dell'editor che
            consentono di applicare stili e struttura al testo in Markdown.
            Nel progetto include: grassetto, corsivo, sottolineato, barrato,
            titoli, link ipertestuali, elenchi puntati e numerati, tabelle e
            citazioni. La gestione delle immagini è esplicitamente esclusa dai
            requisiti.
      \item \textbf{Framework:} Struttura software di base o piattaforma che fornisce librerie e strumenti per facilitare e standardizzare lo sviluppo di applicazioni.
\end{itemize}
 
% ------------- G -------------
\phantomsection
\addcontentsline{toc}{subsection}{G}
\subsection*{G}
\begin{itemize}[leftmargin=*]
      \item \textbf{GitHub Pages:} Servizio di GitHub che consente di pubblicare siti
            web statici direttamente da un repository. Utilizzato dal gruppo per
            esporre la documentazione prodotta come sito web accessibile pubblicamente.
      \item \textbf{GitHub Project:} Strumento di gestione dei task integrato in
            GitHub, utilizzato per pianificare, tracciare e monitorare le attività
            del progetto tramite board e issue.
      \item \textbf{Google Calendar:} Servizio di calendario condiviso di Google,
            utilizzato dal gruppo per la gestione delle disponibilità dei membri e
            la pianificazione delle riunioni ricorrenti e straordinarie.
      \item \textbf{Google Form:} Strumento di Google utilizzato per la creazione
          guidata dei task; genera automaticamente una issue su GitHub e registra
          i dati nel foglio di calcolo condiviso per il monitoraggio delle
          misurazioni.
\end{itemize}
 
% ------------- H -------------
\phantomsection
\addcontentsline{toc}{subsection}{H}
\subsection*{H}
\begin{itemize}[leftmargin=*]
    \item \textbf{HTML:} (HyperText Markup Language) Linguaggio di formattazione utilizzato per strutturare e visualizzare documenti sul World Wide Web.
\end{itemize}
 
% ------------- I -------------
\phantomsection
\addcontentsline{toc}{subsection}{I}
\subsection*{I}
\begin{itemize}[leftmargin=*]
      \item \textbf{Interfaccia Utente (UI):} Insieme degli elementi grafici e interattivi attraverso cui l'utente interagisce con il sistema. Deve garantire semplicità, chiarezza e accessibilità.
      \item \textbf{ISO/IEC 12207:1995:} Standard internazionale
            (\textit{Information Technology~-- Software Life Cycle Processes})
            che definisce un framework di processi per il ciclo di vita del software,
            suddivisi in processi primari (es.\ fornitura, sviluppo), di supporto
            (es.\ documentazione, verifica, gestione della configurazione) e
            organizzativi (es.\ gestione, infrastruttura). Fornisce un quadro
            metodologico per strutturare, gestire e controllare le attività di
            sviluppo e manutenzione del software, applicabile tramite tailoring al
            contesto specifico del progetto.
      \item \textbf{Issue:} Unità di lavoro o segnalazione registrata su GitHub,
            utilizzata per tracciare attività, problemi o richieste. Nel flusso del
            gruppo ogni issue corrisponde a un task e viene generata automaticamente
            tramite Google Form.
      \item \textbf{Issue Tracking:} Processo di monitoraggio delle attività tramite issue GitHub, includendo stati progressivi, assegnatari, tempi previsti ed effettivi, e collegamento con il sistema di misurazione ore.
\end{itemize}
 
% ------------- J -------------
\phantomsection
\addcontentsline{toc}{subsection}{J}
\subsection*{J}
\begin{itemize}[leftmargin=*]
      \item \textbf{JSON:} (JavaScript Object Notation) Formato leggero per lo scambio di dati, facilmente leggibile dalle macchine e dagli esseri umani.
      \item \textbf{Just-in-time:} Principio organizzativo secondo cui ogni attività
          viene normata immediatamente prima di essere attuata, evitando
          pianificazioni eccessive in anticipo. Applicato al \textit{way of working},
          implica che le norme vengano estese e aggiornate in modo incrementale,
          a intervalli brevi, man mano che nuove attività vengono introdotte.
\end{itemize}
 
% ------------- K -------------
\phantomsection
\addcontentsline{toc}{subsection}{K}
\subsection*{K}
\textit{Nessun termine.}
 
% ------------- L -------------
\phantomsection
\addcontentsline{toc}{subsection}{L}
\subsection*{L}
\begin{itemize}[leftmargin=*]
      \item \textbf{LLM (Large Language Model):} Modello avanzato di intelligenza
          artificiale addestrato su grandi quantità di testo per la comprensione,
          elaborazione e generazione di linguaggio umano in modo contestuale.
      \item \textbf{Lettera di presentazione:} Documento formale redatto dal gruppo per accompagnare la candidatura, contenente la motivazione della scelta del capitolato e i collegamenti ai documenti prodotti.
      \item \textbf{Logout:} Operazione con cui un utente termina volontariamente la propria sessione all'interno dell'applicazione, invalidando eventuali token o credenziali attive e impedendo ulteriori interazioni fino a una nuova autenticazione.
      \item \textbf{Lucidspark:} Strumento online per la creazione collaborativa di diagrammi, utilizzato dal gruppo per la modellazione UML e la rappresentazione dei Casi d'Uso.
\end{itemize}
 
% ------------- M -------------
\phantomsection
\addcontentsline{toc}{subsection}{M}
\subsection*{M}
\begin{itemize}[leftmargin=*]
\item \textbf{Merge:} Operazione Git che integra le modifiche di un branch
      di lavoro nel branch principale (\texttt{main}). Nel flusso del gruppo
      il merge viene eseguito dall'Assegnatario del task (non dal Verificatore)
      dopo che la pull request ha ricevuto l'approvazione del Verificatore.
      \item \textbf{Messaggi Precompilati:} Elementi interattivi (es.\ bottoni) che rappresentano azioni o richieste già definite, proposte all'utente per guidare l'interazione e ridurre l'ambiguità del linguaggio naturale.
      \item \textbf{Milestone:} Punto di controllo formale nel calendario di progetto che
      segna il completamento di una fase o il raggiungimento di un obiettivo
      significativo. Nel progetto le milestone principali sono la consegna RTB
      (giugno--luglio 2026) e la consegna PB (agosto 2026).
      \item \textbf{Misurazione:} Raccolta strutturata di dati quantitativi relativi
          allo svolgimento delle attività del gruppo, comprensiva di: ruolo
          ricoperto, tempo previsto e tempo effettivo impiegato. I dati sono
          registrati nel foglio di calcolo condiviso e utilizzati per il
          monitoraggio dell'avanzamento.
      \item \textbf{MongoDB:} Database NoSQL orientato ai documenti, open source,progettato per scalabilità orizzontale e alte prestazioni in contesti distribuiti.
      \item \textbf{MVP (Minimum Viable Product):} Versione del prodotto dotata delle
          funzionalità minime sufficienti per essere valutata dal proponente.
          Consente di verificare la bontà della visione iniziale, di raccogliere
          feedback fondato e di prendere decisioni informate per il completamento
          definitivo del prodotto.
      \item \textbf{Main:} Ramo (branch) principale del repository Git, che contiene esclusivamente le versioni stabili, verificate e approvate dei documenti e del codice.
      \item \textbf{Markdown:} Linguaggio di markup leggero che permette di formattare il testo utilizzando una sintassi semplice e intuitiva in puro testo semplice.
\end{itemize}
 
% ------------- N -------------
\phantomsection
\addcontentsline{toc}{subsection}{N}
\subsection*{N}
\begin{itemize}[leftmargin=*]
    \item \textbf{Norme Di Progetto:} Documento gestionale interno che definisce
          le procedure operative, le convenzioni redazionali e gli strumenti adottati
          dal gruppo, organizzati secondo i processi del ciclo di vita definiti da
          ISO/IEC 12207:1995. Redatto e mantenuto dall'Amministratore; soggetto
          ad aggiornamento incrementale secondo il principio just-in-time.
\end{itemize}
 
% ------------- O -------------
\phantomsection
\addcontentsline{toc}{subsection}{O}
\subsection*{O}
\begin{itemize}[leftmargin=*]
    \item \textbf{OCR (Optical Character Recognition):} Tecnologia che permette di riconoscere e convertire testo contenuto in immagini o documenti scansionati in dati testuali modificabili.
    \item \textbf{OdG Preliminare:} Acronimo di ``Ordine del Giorno Preliminare''.
          Indica l'elenco dei punti da trattare nella riunione, pubblicato dal
          Responsabile prima dell'inizio della seduta su un documento Google
          condiviso. Funge anche da base per la successiva redazione del verbale
          interno, in quanto nel medesimo documento vengono annotate le decisioni
          adottate durante la riunione.
    \item \textbf{OWASP (Open Web Application Security Project):} Fondazione no-profit che produce standard, strumenti e linee guida per la sicurezza delle applicazioni web.
\end{itemize}
 
% ------------- P -------------
\phantomsection
\addcontentsline{toc}{subsection}{P}
\subsection*{P}
\begin{itemize}[leftmargin=*]
      \item \textbf{Pair Programming:} Tecnica di sviluppo software in cui due membri del
      team lavorano insieme sullo stesso task, condividendo lo stesso
      strumento di lavoro: uno scrive (driver) e l'altro revisiona in tempo reale
      (navigator). Utilizzata come misura di mitigazione per i rischi legati
      all'inesperienza tecnologica.
      \item \textbf{Pannello contestuale:} Elemento dell'interfaccia utente che compare in risposta a un'azione dell'utente per mostrare informazioni, suggerimenti o funzionalità rilevanti rispetto al contesto corrente. Nel progetto viene utilizzato per presentare i risultati delle analisi AI o opzioni operative senza interrompere il flusso di lavoro.
      \item \textbf{PB (Product Baseline):} Seconda revisione formale di avanzamento
      del progetto, successiva alla RTB. Fissa il prodotto software completo,
      includendo architettura, codice e documentazione tecnica finale. Nel
      calendario del gruppo è prevista per agosto 2026.
      \item \textbf{Pedice G:} Convenzione tipografica adottata nelle norme di progetto
          per segnalare che un termine è definito nel Glossario. Viene applicato
          alla prima occorrenza del termine in ciascun documento tramite lo script
          Python di automazione (\texttt{python\_glossary\_automation/}).
          \textit{Esempio}: Termine\textsubscript{G}.
      \item \textbf{Piano di Progetto:} Documento gestionale e strategico mirato all'efficienza del team. Definisce l'organigramma, l'assegnazione e la rotazione dei ruoli, il preventivo dell'impegno orario e il consuntivo di periodo. Gestisce l'analisi dei rischi e struttura le attività in una pianificazione macroscopica a lungo termine (fissando le milestone a ritroso) e una pianificazione dettagliata a breve termine (guardando in avanti per i singoli sprint).
      \item \textbf{Piano Di Qualifica:} Documento gestionale esterno che definisce
      gli obiettivi quantitativi di qualità di processo e di prodotto, le
      metriche adottate per misurarli e il cruscotto di valutazione utilizzato
      per monitorarne il raggiungimento nel tempo. Include le retrospettive
      e le iniziative di automiglioramento del gruppo.
      \item \textbf{PoC (Proof of Concept):} Artefatto eseguibile, da considerarsi
          usa-e-getta sul piano architetturale, realizzato all'inizio del progetto
          per valutare la fattibilità tecnologica del prodotto atteso rispetto a
          specifici sottoinsiemi di funzionalità individuati con il proponente.
          Costituisce la technology baseline e aiuta a consolidare l'analisi dei
          requisiti dal lato della soluzione.
      \item \textbf{Pop-up:} Finestra grafica temporanea che appare in sovrapposizione all'interfaccia principale per mostrare informazioni o risultati contestuali senza interrompere il flusso di lavoro dell'utente.
      \item \textbf{Prompt Engineering:} Attività di progettazione e ottimizzazione dei prompt forniti a un LLM per ottenere risposte coerenti, accurate e aderenti agli obiettivi del sistema.
      \item \textbf{Postgres (PostgreSQL):} Sistema avanzato di gestione di database
          relazionali e ad oggetti, open source. Utilizzato come tecnologia di
          persistenza dei dati nel capitolato di riferimento.
      \item \textbf{Preventivo:} Stima delle ore e dei costi per ogni membro e ruolo
            per uno sprint, calcolata a partire dai tempi previsti delle task assegnate.
            Distinto dal \textit{Preventivo Dei Costi} della Dichiarazione degli Impegni,
            che è il preventivo globale di progetto; il preventivo di sprint è prodotto
            automaticamente dal foglio di calcolo condiviso.
      \item \textbf{Preventivo Dei Costi:} Stima dell'impegno economico complessivo
            del progetto, calcolata moltiplicando le ore previste per ciascun ruolo
            per il relativo costo orario ufficiale. Presentato nella Dichiarazione
            degli Impegni; il valore accettato all'aggiudicazione diventa limite
            superiore invalicabile durante lo svolgimento del progetto.
      \item \textbf{Progettista:} Ruolo del gruppo responsabile delle attività di
            progettazione architetturale (\textit{design}). Attivo dalla baseline RTB.
      \item \textbf{Programmatore:} Ruolo del gruppo responsabile delle attività di
            codifica del prodotto software; attivato durante la fase di sviluppo.
      \item \textbf{Prompt:} Input testuale in linguaggio naturale fornito da un
            utente a un sistema di intelligenza artificiale per guidarne la
            generazione di contenuto.
      \item \textbf{Pull Request:} Meccanismo di GitHub utilizzato per proporre
          l'integrazione di modifiche da un branch di lavoro al branch
          \texttt{main}. Nel flusso documentale del gruppo costituisce il punto
          di ingresso obbligatorio per la verifica: il Verificatore esamina le
          modifiche, lascia eventuali commenti e approva o rifiuta l'integrazione.
      \item \textbf{Project Manager:} Figura responsabile della pianificazione, esecuzione e chiusura di un progetto. Nel gruppo, questa funzione è ricoperta dal Responsabile.
\end{itemize}
 
% ------------- Q -------------
\phantomsection
\addcontentsline{toc}{subsection}{Q}
\subsection*{Q}
\textit{Nessun termine.}
 
% ------------- R -------------
\phantomsection
\addcontentsline{toc}{subsection}{R}
\subsection*{R}
\begin{itemize}[leftmargin=*]
      \item \textbf{Redattore:} Ruolo del gruppo responsabile della produzione del
          contenuto di un documento, nel rispetto del template approvato e delle
          convenzioni redazionali definite nelle Norme di Progetto. Non può
          verificare il proprio lavoro.
      \item \textbf{Redo:} Funzionalità dell'editor che permette rispettivamente di ripristinare una modifica annullata. Supporta la gestione incrementale delle revisioni del testo durante la modifica delle note.
      \item \textbf{Refactoring:} Attività di ristrutturazione del codice sorgente
      finalizzata a migliorarne la leggibilità, la manutenibilità e la qualità
      interna, senza modificarne il comportamento esterno. Utilizzata per ridurre
      il debito tecnico accumulato.
      \item \textbf{Repository:} Archivio Git centralizzato che contiene documenti,
          codice sorgente e cronologia completa delle versioni del progetto.
          Ospitato su GitHub; il branch \texttt{main} contiene esclusivamente le
          versioni approvate e stabili.
      \item \textbf{Responsabile:} Ruolo del gruppo responsabile del coordinamento
      delle attività, dell'elaborazione di piani e scadenze, dell'approvazione
      del rilascio di prodotti parziali o finali e della rappresentanza del
      gruppo verso committente e proponente. Redige il verbale interno di ogni
      riunione e il Diario di Bordo.
      \item \textbf{Revisione:} Fase del ciclo di vita documentale in cui il
          Verificatore controlla struttura, coerenza, correttezza linguistica e
          conformità del documento alle Norme di Progetto, prima che esso venga
          pubblicato nel branch \texttt{main}.
      \item \textbf{Risorse Residue:} Ore e budget ancora disponibili per ciascun ruolo
      e per ciascun membro del team, calcolati sottraendo le ore effettivamente
      rendicontate da quelle preventivate. Monitorate al termine di ogni sprint
      tramite le tabelle di prospetto nel Piano di Progetto.
      \item \textbf{RTB (Requirements and Technology Baseline):} Prima revisione
          formale di avanzamento del progetto. Fissa i requisiti da soddisfare
          (Analisi dei Requisiti) e le tecnologie adottate, dimostrate tramite
          Proof of Concept. Si articola in due parti: valutazione del docente
          Cardin e presentazione al docente Vardanega.
      \item \textbf{Ruolo:} Funzione formalmente assegnata a un membro del gruppo
          per lo svolgimento di specifiche attività. I ruoli previsti sono:
          Responsabile, Amministratore, Analista, Progettista, Programmatore,
          Verificatore. Ogni membro deve ricoprire un solo ruolo alla volta e
          ruotare su tutti i ruoli nel corso del progetto.
      \item \textbf{Render (o Rendering):} Processo di conversione e visualizzazione a schermo del codice sorgente (es. testo in Markdown) in un formato grafico leggibile e impaginato per l'utente finale.
\end{itemize}
 
% ------------- S -------------
\phantomsection
\addcontentsline{toc}{subsection}{S}
\subsection*{S}
\begin{itemize}[leftmargin=*]
      \item \textbf{SAL (Stato Avanzamento Lavori):} Momento di verifica periodica o documento che attesta lo stato di completamento delle attività di progetto rispetto alla pianificazione.
      \item \textbf{Scenario Principale:} Sequenza di passi che descrive il flusso standard di un Caso d'Uso, dal punto di vista dell'attore.
      \item \textbf{Scenari Alternativi:} Varianti del flusso principale che descrivono eccezioni, estensioni o condizioni particolari del Caso d'Uso.
      \item \textbf{Scope Creep:} Fenomeno per cui l'ambito di un progetto si espande
      in modo incontrollato durante lo sviluppo, a causa dell'aggiunta di
      funzionalità non pianificate o non approvate formalmente. Costituisce uno
      dei rischi principali del progetto (RR2) e viene contrastato tramite un
      backlog rigoroso e l'approvazione obbligatoria del
      Responsabile per ogni nuova feature.
      \item \textbf{Scrum:} Modello di sviluppo agile basato su iterazioni brevi (sprint), ruoli definiti e revisione incrementale del prodotto. Nel progetto è adottato con rotazione settimanale dei ruoli.
      \item \textbf{Second Brain:} Nome del capitolato d'appalto scelto dal gruppo, proposto dall'azienda Zucchetti.
      \item \textbf{Sei Cappelli per Pensare:} Metodologia ideata da Edward de Bono che struttura il pensiero in sei prospettive distinte (logica, emotiva, critica, creativa, organizzativa e informativa). Nel progetto viene applicata tramite prompt dedicati per analizzare un testo secondo diverse angolazioni.
      \item \textbf{Semantic Versioning:} Convenzione di versionamento adottata dal
          gruppo basata sul formato \texttt{vX.Y.Z}, dove \texttt{X} (Major)
          identifica rilasci ufficiali o cambiamenti strutturali profondi,
          \texttt{Y} (Minor) indica aggiunte o modifiche significative di contenuto
          e \texttt{Z} (Patch) segnala correzioni ortografiche o modifiche non
          semantiche.
      \item \textbf{Single Point of Failure:} Componente o risorsa del progetto la cui
      indisponibilità causa il blocco dell'intera attività correlata. Nel contesto
      del gruppo si riferisce alla situazione in cui un task
      è assegnato a un unico membro senza documentazione condivisa, esponendo il
      team al rischio RO1.
      \item \textbf{Solo Editor:} Modalità che mostra esclusivamente l’area di editing.
      \item \textbf{Solo Render:} Modalità che mostra esclusivamente la versione renderizzata del documento.
      \item \textbf{Spike Tecnico:} Attività esplorativa a tempo limitato svolta dal team
      per valutare la fattibilità o stimare la complessità di una tecnologia o di
      una funzionalità, prima di inserirla nel backlog come task
      pianificabile. Produce conoscenza, non deliverable.
      \item \textbf{Split View:} Modalità dell’editor che mostra simultaneamente la vista di editing e la vista renderizzata del documento.
      \item \textbf{Sprint:} Iterazione di lavoro breve e limitata nel tempo tipica delle metodologie agili, in cui il team si impegna a completare un set predefinito di task.
      \item \textbf{Sprint Backlog:} Sottoinsieme del backlog generale selezionato
      durante lo Sprint Planning, contenente le task che il gruppo si impegna
      a completare entro lo sprint corrente. Definisce l'obiettivo dello sprint
      e viene monitorato durante la Sprint Review.
      \item \textbf{Sprint Planning:} Cerimonia agile svolta all'inizio di ogni
            sprint in cui il gruppo seleziona le attività da completare nel periodo
            e le assegna ai membri, definendo scadenze e priorità.
      \item \textbf{Sprint Retrospective:} Cerimonia agile svolta al termine di
            ogni sprint in cui il gruppo analizza il proprio modo di lavorare,
            identifica criticità e definisce azioni di miglioramento per lo sprint
            successivo.
      \item \textbf{Sprint Review:} Cerimonia agile svolta al termine di ogni sprint
            in cui il gruppo valuta i risultati ottenuti rispetto agli obiettivi
            pianificati, confrontando avanzamento atteso e avanzamento conseguito.
      \item \textbf{Stack Tecnologico:} Insieme delle tecnologie, strumenti e framework utilizzati per sviluppare un sistema software.
      \item \textbf{Stakeholder:} Individuo, gruppo o organizzazione che ha un interesse attivo nel progetto e può influenzarne (o esserne influenzato da) i risultati, come il proponente o il committente.
\end{itemize}
 
% ------------- T -------------
\phantomsection
\addcontentsline{toc}{subsection}{T}
\subsection*{T}
\begin{itemize}[leftmargin=*]
      \item \textbf{Tailoring:} Processo di adattamento di uno standard (nel caso del
            gruppo: ISO/IEC 12207:1995) al contesto specifico del progetto, mediante
            la selezione motivata dei processi da attivare e l'esclusione giustificata
            di quelli non applicabili.
      \item \textbf{Task:} Unità operativa di lavoro assegnata a un membro del
      gruppo con ruolo, scadenza, stima del tempo previsto e tracciamento
      del tempo effettivo. Creata tramite Google Form, genera automaticamente
      una issue su GitHub; il branch corrispondente segue la convenzione
      \texttt{<id>-<gestionale|tecnico>\_<esterno|interno>\_<nomeDoc>}.
      Tracciata su GitHub Project.
      \item \textbf{Technology Baseline:} Baseline che fissa le tecnologie, i
            framework e le librerie adottate per il progetto, dimostrando tramite
            Proof of Concept la loro adeguatezza e interoperabilità rispetto ai
            requisiti del capitolato. Costituisce uno degli obiettivi della
            revisione RTB.
      \item \textbf{Template:} Struttura predefinita e condivisa utilizzata per la
            redazione di tutti i documenti del gruppo, al fine di garantire uniformità
            formale e coerenza tra i diversi elaborati.
      \item \textbf{Tempo Effettivo:} Tempo realmente impiegato da un membro del
            gruppo per completare un task, espresso in minuti. Registrato nel foglio
            di calcolo condiviso a consuntivo dell'attività.
      \item \textbf{Tempo Previsto:} Stima iniziale del tempo necessario per
            completare un task, espressa in minuti. Definita al momento della
            creazione del task e utilizzata come riferimento per il monitoraggio
            dell'avanzamento.
      \item \textbf{Ticket:} Segnalazione o richiesta di lavoro tecnica tracciata all'interno di un sistema di gestione. Nel flusso del gruppo, equivale a una Issue.
      \item \textbf{Timeout LLM:} Errore generato quando il modello linguistico non risponde entro il tempo massimo previsto, gestito tramite messaggi di errore e retry.
      \item \textbf{Title Case:} Convenzione tipografica adottata dal gruppo per i
            titoli di sezioni e sotto-sezioni, che prevede l'iniziale maiuscola per
            ogni parola. \textit{Esempio}: ``Ciclo Di Vita Documentale''.
      \item \textbf{Traduzione:} Funzionalità del sistema che consente di
            convertire il testo presente nell'editor in una lingua diversa
            tramite invocazione di un LLM. Le lingue supportate sono Inglese,
            Francese e Spagnolo come requisito obbligatorio; il Tedesco è
            previsto come estensione facoltativa. Sono escluse le lingue con
            alfabeti non latini.
      \item \textbf{Trunk-based Development:} Strategia di branching in cui ogni
          membro lavora su un branch dedicato a un singolo documento o attività e
          integra le modifiche nel branch principale (\texttt{main}) tramite pull
          request al termine del lavoro, eliminando poi il branch.
\end{itemize}
 
% ------------- U -------------
\phantomsection
\addcontentsline{toc}{subsection}{U}
\subsection*{U}
\begin{itemize}
      \item \textbf{UML (Unified Modeling Language):} Linguaggio di modellazione
      standard utilizzato per rappresentare graficamente la struttura e il
      comportamento di un sistema software. Nel contesto del progetto è
      adottato per la rappresentazione dei diagrammi dei casi d'uso.
      \item \textbf{Undo:} Funzionalità dell'editor che permette rispettivamente di annullare l'ultima modifica effettuata. Supporta la gestione incrementale delle revisioni del testo durante la modifica delle note.
      \item \textbf{Usabilità:} Capacità di un sistema di essere utilizzato in modo efficace, efficiente e soddisfacente dall'utente finale.
      \item \textbf{Utenza Target:} Profilo dell'utente finale a cui il sistema
      è destinato, definito in accordo con il proponente. Nel contesto del
      progetto Second Brain, l'utente target è un utente giovane,
      tipicamente uno studente, che utilizza l'applicazione per prendere
      appunti e desidera essere supportato nella loro elaborazione tramite
      LLM.
\end{itemize}
 
% ------------- V -------------
\phantomsection
\addcontentsline{toc}{subsection}{V}
\subsection*{V}
\begin{itemize}[leftmargin=*]
      \item \textbf{Valutazione Dei Capitolati:} Documento gestionale esterno prodotto
          in fase di candidatura che analizza e confronta i capitolati proposti
          dalle aziende proponenti, includendo il resoconto degli incontri
          esplorativi, i criteri di valutazione, i vantaggi e gli svantaggi di
          ciascuna proposta e la motivazione della scelta finale.
      \item \textbf{VarGroup:} Azienda proponente di uno dei capitolati d'appalto analizzati dal gruppo durante la fase di candidatura.
      \item \textbf{VE-y.x:} Codice identificativo univoco di una decisione adottata
          in un incontro esterno (con proponente o docente), registrata nel relativo
          verbale esterno. \texttt{y} è il numero progressivo del verbale esterno
          e \texttt{x} il numero progressivo della decisione al suo interno.
      \item \textbf{Velocity:} Misura della quantità di lavoro completata dal team in uno
      sprint, espressa in ore o story point. Viene monitorata sprint
      per sprint per rilevare andamenti anomali e per affinare le stime dei sprint
      successivi.
      \item \textbf{Verificatore:} Ruolo del gruppo responsabile del controllo di
          conformità di documenti e codice prodotti dagli altri membri. Deve essere
          sempre un membro diverso dal Redattore dell'elaborato in esame; non può
          verificare il proprio lavoro.
      \item \textbf{VI-y.x:} Codice identificativo univoco di una decisione adottata
          in una riunione interna, registrata nel relativo verbale interno.
          \texttt{y} è il numero progressivo del verbale interno e \texttt{x} il
          numero progressivo della decisione al suo interno.
\end{itemize}
 
% ------------- W -------------
\phantomsection
\addcontentsline{toc}{subsection}{W}
\subsection*{W}
\begin{itemize}[leftmargin=*]
    \item \textbf{Way of Working (WoW):} Insieme strutturato e documentato di
          regole, strumenti, procedure e metodologie adottate dal gruppo per
          organizzare e svolgere il lavoro in modo uniforme e tracciabile.
          Corrisponde operativamente alle Norme di Progetto e si evolve in modo
          incrementale secondo il principio just-in-time per tutta la durata del
          progetto.
    \item \textbf{WhatsApp:} Applicazione di messaggistica istantanea utilizzata
          dal gruppo come canale secondario per comunicazioni rapide, notifiche
          urgenti e convocazioni di riunioni straordinarie.
\end{itemize}
 
% ------------- X -------------
\phantomsection
\addcontentsline{toc}{subsection}{X}
\subsection*{X}
\textit{Nessun termine.}
 
% ------------- Y -------------
\phantomsection
\addcontentsline{toc}{subsection}{Y}
\subsection*{Y}
\textit{Nessun termine.}
 
% ------------- Z -------------
\phantomsection
\addcontentsline{toc}{subsection}{Z}
\subsection*{Z}
\begin{itemize}[leftmargin=*]
    \item \textbf{Zucchetti:} Azienda proponente del capitolato "Second Brain", per il quale il gruppo ha scelto di presentare la propria candidatura.
\end{itemize}